% ============================================================
% MÓDULO: FORMAS JURÍDICAS DEL EMPRENDIMIENTO
% ============================================================
\clearpage
\section*{MÓDULO: FORMAS JURÍDICAS DEL EMPRENDIMIENTO}
\subsection*{Monotributo, autónomo, sociedades y cooperativas}

\begin{recuadro}
\textbf{Objetivo del módulo:} conocer las distintas formas legales con las que se puede llevar
adelante un emprendimiento, cuánto cuesta y cuánto tiempo lleva cada una, qué
\textbf{responsabilidad} asumen sus dueños y cómo \textbf{tributa} cada una, para elegir la
más conveniente según la etapa y el tamaño del proyecto.
\end{recuadro}

% ------------------------------------------------------------
\section*{1. DOS PREGUNTAS CLAVE}

Antes de elegir una forma jurídica conviene responder:
\begin{enumerate}
  \item \textbf{¿Emprendo solo o con socios?}
  \item \textbf{¿Qué pasa con mis bienes personales si al negocio le va mal?} Esto se llama
        \textbf{responsabilidad}:
  \begin{itemize}
    \item \textbf{Ilimitada:} las deudas del negocio se pueden cobrar con los bienes
          personales del dueño (su auto, su casa, sus ahorros).
    \item \textbf{Limitada:} los socios arriesgan solo el capital que aportaron. Las deudas
          de la sociedad las paga la sociedad con su patrimonio.
  \end{itemize}
\end{enumerate}

Una \textbf{persona humana} (antes llamada ``persona física'') trabaja a su propio nombre. Una
\textbf{persona jurídica} (una sociedad o una cooperativa) es un sujeto distinto de sus socios:
tiene su propio CUIT, sus propios bienes, sus propias deudas y sus propios impuestos.

% ------------------------------------------------------------
\section*{2. PERSONA HUMANA}

\subsubsection*{2.1. Monotributista}
Ya estudiado en la unidad de Sistema Impositivo. Es el régimen simplificado para pequeños
contribuyentes: una \textbf{cuota mensual única} que reúne el componente impositivo (en lugar de
IVA y Ganancias), la jubilación y la obra social, según la categoría por facturación.

\begin{center}
\small
\begin{tabular}{|p{0.44\linewidth}|p{0.44\linewidth}|}
\hline
\textbf{Ventajas} & \textbf{Desventajas} \\ \hline
\begin{itemize}[leftmargin=*,nosep]
  \item Alta inmediata y gratuita en ARCA.
  \item Cuota fija y previsible; no se liquida IVA ni Ganancias.
  \item Incluye obra social y aportes jubilatorios.
  \item Pocas obligaciones formales: facturar y recategorizarse.
\end{itemize}
&
\begin{itemize}[leftmargin=*,nosep]
  \item Topes de facturación: si se superan, se pasa a Responsable Inscripto.
  \item Responsabilidad \textbf{ilimitada}.
  \item No discrimina IVA: a un cliente Responsable Inscripto le resulta menos conveniente.
  \item Jubilación mínima.
\end{itemize}
\\ \hline
\end{tabular}
\end{center}

\subsubsection*{2.2. Autónomo / Responsable Inscripto}
Cuando se superan los topes del monotributo (o por elección), la persona humana se inscribe
como \textbf{Responsable Inscripto} en IVA y Ganancias y como \textbf{trabajador autónomo} en
el sistema previsional.
\begin{itemize}
  \item \textbf{IVA:} se cobra a los clientes (21\% general, 10,5\% reducido) y se descuenta lo
        pagado a proveedores. Declaración jurada mensual.
  \item \textbf{Ganancias:} escala progresiva para personas humanas (del 5\% al 35\%), con
        deducciones personales. Declaración jurada anual y anticipos.
  \item \textbf{Aportes de autónomos:} cuota mensual según la categoría de autónomo, más la
        obra social que se elija.
  \item \textbf{Ingresos Brutos:} régimen general provincial (en la Provincia de Buenos Aires,
        ARBA).
\end{itemize}

\begin{center}
\small
\begin{tabular}{|p{0.44\linewidth}|p{0.44\linewidth}|}
\hline
\textbf{Ventajas} & \textbf{Desventajas} \\ \hline
\begin{itemize}[leftmargin=*,nosep]
  \item Sin tope de facturación.
  \item Emite factura A: el cliente empresa computa el IVA.
  \item Puede descontar el IVA de compras y deducir gastos en Ganancias.
\end{itemize}
&
\begin{itemize}[leftmargin=*,nosep]
  \item Carga impositiva y administrativa mucho mayor: necesita contador.
  \item Responsabilidad \textbf{ilimitada}.
  \item Ganancias con alícuotas de hasta el 35\%.
\end{itemize}
\\ \hline
\end{tabular}
\end{center}

% ------------------------------------------------------------
\section*{3. SOCIEDADES}

La \textbf{Ley General de Sociedades N° 19.550} regula los tipos societarios. Una sociedad
nace de un \textbf{contrato} (o \textbf{estatuto}) que define socios, capital, objeto (a qué se
dedica), administración y forma de repartir las ganancias. Las sociedades regulares se
inscriben en el registro público: la \textbf{IGJ} en la Ciudad de Buenos Aires o la
\textbf{DPPJ} (Dirección Provincial de Personas Jurídicas) en la Provincia de Buenos Aires.

\subsubsection*{3.1. Sociedad de hecho (Sociedades de la Sección IV)}
Antes de 2015 se llamaba \textbf{sociedad de hecho} a la que funcionaba sin contrato escrito:
dos o más personas que trabajan juntas y comparten gastos y ganancias. Con el Código Civil y
Comercial, hoy se las denomina \textbf{sociedades de la Sección IV} (también ``simples'' o
``informales''): las que no adoptan un tipo de la ley, omiten requisitos o no se inscriben.

\begin{itemize}
  \item Los socios responden con sus bienes personales por las deudas sociales, en general
        \textbf{en partes iguales} (mancomunadamente), salvo que se haya pactado otra cosa.
  \item Pueden adherir al \textbf{monotributo} con un máximo de \textbf{3 socios}: la sociedad
        paga el componente impositivo incrementado un 20\% por cada socio, y cada socio paga
        sus propios aportes jubilatorios y de obra social.
  \item Si superan los topes, tributan como Responsable Inscripto (IVA) y Ganancias.
\end{itemize}
\textbf{Útil para:} dos o tres socios que empiezan, con poca inversión y poco riesgo.
\textbf{Riesgo:} sin contrato claro, los conflictos entre socios (quién aportó qué, cómo se
reparte, cómo se sale) son difíciles de resolver.

\subsubsection*{3.2. S.R.L. (Sociedad de Responsabilidad Limitada)}
\begin{itemize}
  \item Entre \textbf{2 y 50 socios}. El capital se divide en \textbf{cuotas}.
  \item Responsabilidad \textbf{limitada} al capital aportado.
  \item Administración a cargo de uno o más \textbf{gerentes}.
  \item Se constituye por instrumento público o privado con firmas certificadas y se inscribe
        en el registro público.
\end{itemize}
\textbf{Útil para:} pymes familiares o con pocos socios que se conocen; talleres, comercios.
\textbf{Contras:} trámite más lento y costoso que una SAS; la transferencia de cuotas a
terceros suele estar restringida.

\subsubsection*{3.3. S.A. (Sociedad Anónima)}
\begin{itemize}
  \item Capital dividido en \textbf{acciones}, que se pueden transferir con facilidad.
  \item Capital mínimo de \textbf{\$30.000.000} (Decreto 209/2024).
  \item Órganos: \textbf{asamblea} de accionistas, \textbf{directorio} y, en ciertos casos,
        \textbf{sindicatura}. Mayores controles del registro público.
  \item Existe la \textbf{S.A.U.} (unipersonal), con exigencias adicionales.
\end{itemize}
\textbf{Útil para:} empresas medianas y grandes, que buscan inversores o cotizar en bolsa
(aerolíneas, fabricantes, grandes talleres).
\textbf{Contras:} la forma más costosa y burocrática.

\subsubsection*{3.4. S.A.S. (Sociedad por Acciones Simplificada)}
Creada por la \textbf{Ley 27.349 de Apoyo al Capital Emprendedor} (2017), pensada para
emprendedores.
\begin{itemize}
  \item Puede tener \textbf{un solo socio} y un administrador unipersonal.
  \item Capital mínimo: \textbf{dos salarios mínimos, vitales y móviles}.
  \item Constitución \textbf{digital}, con estatuto modelo, en muy pocos días, y CUIT online.
  \item Objeto amplio: en la Ciudad de Buenos Aires, la IGJ admite desde 2026 la fórmula
        ``la realización de cualquier actividad lícita''.
  \item Responsabilidad \textbf{limitada}.
\end{itemize}
\textbf{Útil para:} emprendedores individuales o equipos chicos que quieren separar su
patrimonio del negocio, startups, proyectos que pueden recibir inversores.
\textbf{Contras:} tributa como cualquier sociedad (no puede ser monotributista); lleva
contabilidad, balances y libros digitales; el régimen registral depende de cada jurisdicción.

\subsubsection*{3.5. Cómo tributan las sociedades (SRL, SA, SAS)}
\begin{itemize}
  \item \textbf{IVA:} Responsable Inscripto, igual que un autónomo.
  \item \textbf{Ganancias:} escala propia para sociedades, por tramos de ganancia neta. Para los
        ejercicios iniciados desde el 1/1/2026:
\end{itemize}

\begin{center}
\small
\begin{tabular}{|p{0.50\linewidth}|c|}
\hline
\textbf{Ganancia neta imponible anual} & \textbf{Alícuota marginal} \\ \hline
Hasta \$133.200.243,59 & 25\% \\ \hline
De \$133.200.243,59 a \$1.332.002.435,93 & 30\% \\ \hline
Más de \$1.332.002.435,93 & 35\% \\ \hline
\end{tabular}
\end{center}

\begin{itemize}
  \item \textbf{Dividendos:} cuando la sociedad reparte utilidades a los socios, se retiene un
        \textbf{7\%} adicional.
  \item \textbf{Ingresos Brutos} provincial y, si los socios cobran por dirigir la sociedad,
        deben inscribirse como autónomos.
  \item Los montos de la escala se actualizan todos los años: consultar ARCA.
\end{itemize}

\begin{recuadro}
\textbf{¿Cuándo conviene pasar de monotributo a sociedad?} No hay un número mágico: depende de
la facturación, los márgenes, la cantidad de socios y el riesgo de la actividad. Una regla
práctica es mantenerse como monotributista mientras se valida el negocio y armar una
\textbf{SAS o SRL} cuando aparecen socios, inversiones, clientes grandes que exigen factura A
o una actividad con riesgo de deudas o reclamos importantes. Siempre consultar a un
\textbf{contador público}.
\end{recuadro}

% ------------------------------------------------------------
\section*{4. COOPERATIVAS}

Una \textbf{cooperativa} es una asociación de personas que se unen voluntariamente para
satisfacer necesidades comunes mediante una empresa de propiedad conjunta y gestión
democrática. Se rigen por la \textbf{Ley 20.337} y son registradas y controladas por el
\textbf{INAES} (Instituto Nacional de Asociativismo y Economía Social), que les otorga la
\textbf{matrícula}.

\subsubsection*{4.1. Principios cooperativos}
\begin{enumerate}
  \item Adhesión voluntaria y abierta.
  \item Gestión democrática: \textbf{un asociado, un voto}, sin importar el capital aportado.
  \item Participación económica de los asociados.
  \item Autonomía e independencia.
  \item Educación, capacitación e información.
  \item Cooperación entre cooperativas.
  \item Compromiso con la comunidad.
\end{enumerate}

\subsubsection*{4.2. Tipos de cooperativas}
\begin{center}
\small
\begin{tabular}{|p{0.22\linewidth}|p{0.68\linewidth}|}
\hline
\textbf{Tipo} & \textbf{Qué hace} \\ \hline
De trabajo & Los asociados son los trabajadores: producen bienes o servicios y se reparten
los resultados (talleres, fábricas recuperadas, servicios técnicos). \\ \hline
De provisión & Compran en conjunto insumos o servicios para sus asociados, que son
productores o comerciantes independientes. \\ \hline
De consumo & Compran bienes para venderlos a sus asociados a mejor precio. \\ \hline
De servicios públicos & Electricidad, agua, telefonía, internet en pueblos y ciudades. \\ \hline
Agropecuarias & Productores rurales que acopian, procesan y comercializan juntos. \\ \hline
De vivienda & Construyen o adquieren viviendas para sus asociados. \\ \hline
De crédito y seguros & Préstamos y seguros entre asociados. \\ \hline
Escolares & Formadas por estudiantes, con fines educativos (ver 4.4). \\ \hline
\end{tabular}
\end{center}

\subsubsection*{4.3. Tributación, ventajas y desventajas}
\begin{itemize}
  \item Las cooperativas están \textbf{exentas del Impuesto a las Ganancias}. Los
        \textbf{excedentes} (no se llaman ganancias) se distribuyen entre los asociados en
        proporción al trabajo o uso de los servicios, no al capital.
  \item En \textbf{IVA} e \textbf{Ingresos Brutos} tributan según la actividad (hay exenciones
        o tratamientos especiales para algunas).
  \item En las cooperativas de trabajo los asociados no son empleados: cobran \textbf{retornos}
        o anticipos y se inscriben como \textbf{monotributistas} o como autónomos.
  \item Acceso a programas de promoción y créditos específicos (INAES, provincia).
\end{itemize}
\textbf{Contras:} trámite de constitución más largo (cursos, asamblea, matrícula del INAES),
número mínimo de asociados, controles y balances anuales, decisiones por asamblea y menor
atractivo para inversores, ya que el capital no da más votos.

\subsubsection*{4.4. Cooperativas escolares}
Son cooperativas formadas por \textbf{estudiantes}, con la guía de docentes, cuyo objetivo
principal es \textbf{educativo}: aprender a trabajar en equipo, gestionar un emprendimiento y
practicar la democracia. Por tratarse de menores de 18 años, se rigen por las disposiciones de
la \textbf{autoridad educativa} (en la Provincia de Buenos Aires, la Dirección General de
Cultura y Educación), de acuerdo con los principios de la Ley 20.337.

\begin{itemize}
  \item \textbf{Ejemplos:} venta de útiles y fotocopias, kiosco saludable, huerta, reparación de
        equipos electrónicos para la comunidad, producción en el taller.
  \item \textbf{Órganos:} asamblea de estudiantes asociados, consejo de administración y
        síndico elegidos por los alumnos, con un docente asesor.
  \item No reparten ganancias como una empresa: los excedentes se reinvierten en la escuela o
        en proyectos del grupo.
  \item No confundir con la \textbf{asociación cooperadora}, formada por familias y vecinos
        para apoyar a la escuela.
\end{itemize}

% ------------------------------------------------------------
\section*{5. CUADRO COMPARATIVO}

\begin{center}
\footnotesize
\setlength{\tabcolsep}{3pt}
\begin{tabular}{|>{\raggedright\arraybackslash}p{0.13\linewidth}|>{\raggedright\arraybackslash}p{0.10\linewidth}|>{\raggedright\arraybackslash}p{0.13\linewidth}|>{\raggedright\arraybackslash}p{0.15\linewidth}|>{\raggedright\arraybackslash}p{0.21\linewidth}|>{\raggedright\arraybackslash}p{0.17\linewidth}|}
\hline
\textbf{Forma} & \textbf{Socios} & \textbf{Responsab.} & \textbf{Constitución} &
\textbf{Impuestos} & \textbf{Ideal para} \\ \hline
Monotributo & 1 & Ilimitada & Alta online, gratis & Cuota mensual única & Empezar,
validar la idea \\ \hline
Autónomo (RI) & 1 & Ilimitada & Alta online & IVA, Ganancias 5--35\%, autónomos, IIBB &
Facturación alta sin socios \\ \hline
Sección IV (ex de hecho) & 2 o más (monotributo: hasta 3) & Partes iguales, con bienes
personales & Sin trámite registral & Monotributo o IVA y Ganancias & Socios que recién
empiezan \\ \hline
S.R.L. & 2 a 50 & Limitada & Contrato e inscripción & IVA, Ganancias 25--35\%, IIBB & Pymes
con pocos socios \\ \hline
S.A. & 1 (SAU) o más & Limitada & Escritura e inscripción; capital \$30M & IVA, Ganancias
25--35\%, IIBB & Empresas medianas y grandes \\ \hline
S.A.S. & 1 o más & Limitada & Digital, rápida; capital 2 SMVM & IVA, Ganancias 25--35\%, IIBB
& Emprendedores y startups \\ \hline
Cooperativa & Mínimo legal & Limitada & Matrícula INAES & Exenta de Ganancias; IVA e IIBB
según actividad & Trabajo asociado, servicios a la comunidad \\ \hline
\end{tabular}
\end{center}

% ------------------------------------------------------------
\section*{6. ACTIVIDADES DEL MÓDULO}

\begin{enumerate}
  \item \textbf{¿Qué forma le conviene?} Para cada caso, elegir la forma jurídica y
        justificar:
  \begin{enumerate}
    \item Un técnico que repara radios de aviones de aeroclubes, solo, con poca facturación.
    \item Tres egresados que arman un taller de reparación de electrónica.
    \item Un emprendedor que desarrolla una aplicación y busca inversores.
    \item Doce trabajadores de una empresa que cerró y quieren seguir produciendo.
    \item Un taller aeronáutico con 40 empleados que factura a aerolíneas.
  \end{enumerate}
  \item \textbf{Cálculo de Ganancias:} con la escala de sociedades, calcular el impuesto de una
        SAS que obtuvo una ganancia neta de \$150.000.000 y cuánto recibe el socio si reparte
        toda la utilidad (dividendos al 7\%).
  \item \textbf{Proyecto final:} definir la forma jurídica que adoptaría el emprendimiento del
        grupo, con sus ventajas, desventajas y costos.
  \item \textbf{Cooperativa escolar:} redactar la propuesta de una cooperativa escolar del
        curso (objeto, asociados, órganos, reglas para repartir tareas y reinvertir
        excedentes), por ejemplo un servicio de reparación de equipos electrónicos para la
        comunidad.
\end{enumerate}
